\begin{papercard}{martin2023}
\cf{Stages} \stagetag{S} (what counts as a privacy-enhancing claim; a ``sound conception of privacy'' as the benchmark), \stagetag{H} (inferences and purposes as sources of privacy violation), \stagetag{C} (a vignette design built on the contextual-integrity parameters). A 2023 \emph{Georgetown Law Technology Review} article by Martin, Nissenbaum and Shmatikov that argues, with philosophical analysis plus six factorial-vignette surveys, that three industry ``privacy-enhancing'' principles (restricting third parties, using inferences instead of raw data, keeping data on the device) do not match users' privacy expectations. For the thesis it is chiefly (i) a worked example of studying norms with contextual-integrity vignettes and (ii) evidence that purpose and inference matter more than data type or party.

\cf{Core argument} The authors accept that the techniques may achieve their stated technical goals, but ``the principles themselves fall short when the privacy-enhancing technique does not address privacy-violating behavior'' \pc{martin2023}{1}{we argue that the principles themselves fall short when the privacy-enhancing technique does not address privacy-violating behavior}. Calling a system privacy-enhancing requires that it demonstrably achieves certain flow patterns \emph{and} that those patterns embody a sound conception of privacy; otherwise it is like a perfectly aligned gun sight pointed at the wrong target \pc{martin2023}{5}{Claiming that a given system is privacy-enhancing requires both that it demonstrably achieves certain flow patterns and that the patterns embody a sound conception of privacy}. Their conception is contextual integrity (CI): ``privacy is the appropriate flow of information'' \pc{martin2023}{5}{contextual integrity (CI), basically asserting that privacy is the appropriate flow of information}. Concluding statement: efforts must start with ``the explicit adoption of a meaningful conception of privacy'' \pc{martin2023}{1}{any effort to create successful privacy-enhancing systems must start with the explicit adoption of a meaningful conception of privacy}.

\cf{Three benchmarks for a conception of privacy} Precision (clear, rigorous, internally consistent), ethics (explains privacy's ethical significance), fidelity (roughly faithful to common usage) \pc{martin2023}{6}{These benchmarks are precision (is the conception clear, rigorous, and internally consistent), ethics (does the conception explain privacy’s ethical significance), and fidelity (is the conception roughly faithful to common usage)}; borrowed from analytic philosophy's evaluation of definitions. The empirical studies test the \emph{fidelity} benchmark.

\cf{Three principles examined (Part I, summary)} (1) Limiting access to personal information by third parties; the first-party/third-party distinction is argued to lack moral clarity (a first party such as an ad network can be a non-contextual recipient; third parties such as payment processors and delivery drivers can be contextual) and law and regulation often assume it. (2) Minimising use and retention of raw data in favour of inferences; inferences can be at least as revealing and are governed by weaker rules. (3) Leaving data on the device; for example Google's Privacy Sandbox and federated approaches still enable inference-based targeting.

\cf{Study design: factorial vignettes (Part II)} Method: Jasso's factorial-survey approach. Respondents see 20--40 vignettes in which factors vary systematically and rate ``the degree to which a scenario is appropriate or `okay' '' \pc{martin2023}{34}{After seeing each vignette, respondents are asked to complete a simple rating task—the degree to which a scenario is appropriate or “okay”}. Factors correspond to CI parameters: \emph{contextual actor} (search engine, browser, two social networks, news site), \emph{data type} (search terms, location, likes/clicks, profile information, web activity), \emph{inference} (none; demographics, emotions, friends and activities, interests, medical inferences) and \emph{purpose/use} (improve services, place ads on the site, let others place ads later, sell to others) \pc{martin2023}{36}{Purpose/Use: The purpose of the data collection or derived inference varied from improving services and placing ads on the site to allowing others to later place ads when online and selling the knowledge to others}. Six surveys vary whether inferences are used (not in Surveys 1 and 6; in Surveys 2, 3 and 4) and whether a primary actor or a third party such as a data broker creates and uses them; a seventh survey tests the ``data stays on your device'' design (Survey 5 is the third-party tracker/data-broker comparison) \pc{martin2023}{37}{We ran six studies. In each, the general template was similar and the factors/values were consistent}. Run March--April 2023 on Prolific; each respondent rated thirty vignettes; unique respondents per condition; paid 1.60 GBP (11.50 GBP/h) \pc{martin2023}{38}{The surveys were run in March and April of 2023 using Prolific}. Analysis: regression of the rating on factor levels and t-tests between conditions (Table 2).

\cf{Results} (1) Overall ratings are negative: people do not find the described collection and use ``OK'' on average \pc{martin2023}{39}{In general, the average rating task for all conditions is negative, meaning that individuals do not find the vignettes describing the collection and use of their data to be OK on average}. (2) \emph{Purpose beats data type}: ``changing how knowledge is used had a greater impact than changing the type of information collected, all else being equal'' \pc{martin2023}{39}{changing how knowledge is used had a greater impact than changing the type of information collected, all else being equal}; services-improvement flows were consistently acceptable, whereas non-contextual uses (ads later online, selling) were rated worse \pc{martin2023}{43}{the use of the same information to promote non-contextual purposes, such as to place ads when subsequently online or to sell to others, did not meet privacy expectations}. (3) \emph{Inferences are not a solution}: the mean rating with inferences (Survey 2: $-30.88$) is significantly lower than with raw data (Survey 1: $-23.41$); the pattern holds per data type; medical inferences are worst except that sale to employers or lenders is $-60$ to $-80$ for everything \pc{martin2023}{41}{is significantly lower than the average rating for Survey 1}. (4) \emph{First vs.\ third party}: creation and use of inferences by a data broker is judged as violating as by the first party, and respondents judged a first party \emph{selling} inferences worse than a broker doing so \pc{martin2023}{61}{respondents judged a first party selling their inferences to be worse than the data broker selling their inferences}. (5) \emph{On-device}: a browser that collects data, keeps it locally and creates inferences for ads is rated the same as (mean $-32.0$ vs.\ $-32.7$) third-party trackers feeding a data broker \pc{martin2023}{58}{keep the data on the local device, create inferences for use in advertising}.

\cf{Implications (Part VII)} Practice: ``Simply removing access to raw data—while exposing inferences—is not sufficient to argue that privacy has been protected''; protection depends on which inferences are available (e.g., curation of the Topics API) \pc{martin2023}{63}{Simply removing access to raw data—while exposing inferences—is not sufficient to argue that privacy has been protected}. Policy: rules requiring notice or opt-in only for third-party sharing would not necessarily protect privacy because first parties could still collect and use data for non-contextual purposes \pc{martin2023}{63}{rules or regulations that require notification if sharing with a third party, or require users to opt in to sharing data with a third party, or even limit whether a third party can receive consumer data would not necessarily provide privacy protections}; minimisation and deletion rules should cover inferences \pc{martin2023}{64}{Our findings would support policies that treat inferences just as seriously as underlying data}; purpose was ``even more important ... than what data was collected or with whom the data was shared'' \pc{martin2023}{64}{how information was used, specifically, whether in service of contextual versus non-contextual ends was even more important for our respondents than what data was collected or with whom the data was shared}.

\cf{Limits (card)} Hypothetical vignettes rated by Prolific respondents (convenience sample, English-speaking, demographics not discussed in the extracted text); contexts are search, browsers and social media, not health apps (only one inference type is medical); a single ``OK'' rating per vignette, no behavioural validation; regression results are reported mostly as means and t-tests; the normative argument depends on the CI framework, which not all privacy scholars accept; authors have a stated theoretical commitment (Nissenbaum is the proponent of CI) and the design is built to test CI predictions; US law focus in Part I.

\cf{Use for this thesis} (1) A model for describing information flows in stage H/S by CI parameters (sender/actor, data type, inference, purpose, recipient, transmission principle) rather than by data type alone; this supports coding \emph{purpose} and \emph{inference} for each flow found in policies and screens. (2) Evidence for ranking harms: non-contextual purposes (advertising, resale) and inference-based use are consistently rated worse than first-party service use, and first-party status does not make a flow acceptable. (3) A caution when reading app-store privacy labels: a label claiming ``data not linked'' or ``processed on device'' does not address inference-based profiling. (4) The vignette method is a template if the thesis later adds a small user-judgement component; no such component is assumed. (5) Use together with Nissenbaum (2004, 2009), Doan et al.\ (CI applied to policy text) and Citron \& Solove (harm types).
\end{papercard}
